Los resultados obtenidos muestran una relación aproximadamente lineal entre la
frecuencia de emisión de los LEDs y el voltaje ánodo--cátodo medido en el umbral
de conducción. El ajuste lineal presentó un coeficiente de determinación
$R^2=0.99089$, indicando que el modelo empleado describe adecuadamente el
comportamiento experimental.\\

A partir de la pendiente del ajuste se obtuvo

\[
\left(\frac{h}{e}\right)_{\mathrm{exp}}
=(3.992\pm0.059)\times10^{-15}\ \mathrm{V\,s},
\]

valor que difiere únicamente un $3.47\%$ del valor aceptado
$4.135668\times10^{-15}\ \mathrm{V\,s}$. Esta diferencia resulta razonable para
un experimento de laboratorio y confirma que el procedimiento permite estimar
adecuadamente la relación entre ambas constantes fundamentales.\\

Las principales fuentes de incertidumbre provienen de la determinación del
voltaje umbral de cada LED, la dispersión de las longitudes de onda nominales
proporcionadas por el fabricante, la resolución de los instrumentos de medición
y la aproximación de considerar que toda la energía eléctrica suministrada se
transforma en energía luminosa. En realidad, parte de la energía se disipa como
calor dentro del semiconductor, lo cual introduce una desviación sistemática
respecto al modelo ideal.\\

En el caso del LED infrarrojo se obtuvo un máximo de intensidad en
$\lambda_{\max}=591.5\,\mathrm{nm}$, valor que corresponde a la región
visible del espectro y no a la emisión infrarroja esperada. Esta
discrepancia puede atribuirse a las limitaciones del intervalo espectral
y de la sensibilidad del espectrómetro. Si la longitud de onda nominal
del LED se encuentra alrededor de $940\,\mathrm{nm}$, su máximo de
emisión queda fuera del intervalo registrado, cuyo límite superior fue
aproximadamente $904\,\mathrm{nm}$. Además, la sensibilidad del detector
suele disminuir cerca de los extremos de su intervalo de operación. En
consecuencia, al no registrarse el pico principal del LED, el algoritmo
de análisis identificó como máximo una señal secundaria en la región
visible, posiblemente debida a luz ambiental, ruido instrumental o
emisión espuria. Por esta razón, el valor obtenido no se consideró una
medición representativa de la longitud de onda del LED infrarrojo.\\




